\documentclass[11pt,a4paper,oneside,openright]{book}
\usepackage[spanish,activeacute,es-noshorthands]{babel}
\usepackage[utf8]{inputenc}
\usepackage[T1]{fontenc}
\usepackage{lmodern}
\usepackage{amsmath, amssymb}
\usepackage{graphicx}
\usepackage[dvipsnames]{xcolor}
\usepackage{listings}
\usepackage{geometry}
\geometry{margin=2.5cm, headheight=14pt}
\usepackage{microtype}
\usepackage{booktabs}
\usepackage{enumitem}
\usepackage{fancyhdr}
\usepackage{tcolorbox}
\tcbuselibrary{skins, breakable}
\usepackage{caption}
\usepackage{hyperref}
\hypersetup{colorlinks=true, linkcolor=NavyBlue, urlcolor=BrickRed,
  pdftitle={M5 --- Flight Software para CubeSat}, pdfauthor={INTISAT}}

\definecolor{codebg}{RGB}{248, 249, 250}
\lstdefinestyle{shell}{
  basicstyle=\ttfamily\footnotesize,
  backgroundcolor=\color{codebg},
  frame=single, framesep=4pt, breaklines=true,
  showstringspaces=false,
  literate={a}{{a}}1 {e}{{e}}1 {i}{{i}}1 {o}{{o}}1 {u}{{u}}1 {n}{{n}}1
}
\lstset{style=shell}

\newtcolorbox{intuibox}[1][]{colback=green!5, colframe=green!40!black,
  fonttitle=\bfseries, title=Intuicion, breakable, sharp corners=south, #1}
\newtcolorbox{payload}[1][]{colback=blue!4, colframe=blue!50!black,
  fonttitle=\bfseries, title=Conexion con el INTISAT, breakable, sharp corners=south, #1}
\newtcolorbox{cuidado}[1][]{colback=red!5, colframe=red!50!black,
  fonttitle=\bfseries, title=Cuidado, breakable, sharp corners=south, #1}

\pagestyle{fancy}
\fancyhf{}
\fancyhead[LE,RO]{\thepage}
\fancyhead[RE]{\nouppercase{\leftmark}}
\fancyhead[LO]{\nouppercase{\rightmark}}

\title{\Huge\bfseries M5 --- Flight Software\\[0.4em]
\Large Software de a bordo para CubeSat\\[0.6em]
\large Manual de la coleccion CubeSat (M5 / 11)}
\date{\today}

\begin{document}
\frontmatter
\maketitle

\chapter*{Prefacio}
\addcontentsline{toc}{chapter}{Prefacio}

El \emph{flight software} (FSW) es el software que corre en el OBC del
satelite y en cada subsistema con procesador. Este manual cubre como
diseñarlo, estructurarlo, y operarlo.

\textbf{Prerequisitos}: M0 (CubeSat 101), Raspberry\_Pi\_5\_Profundo.

\tableofcontents

\mainmatter

\chapter{Arquitecturas de OBC}

\section{Microcontrolador vs SoC con Linux}

\begin{center}
\begin{tabular}{p{4cm}p{5cm}p{5cm}}
\toprule
\textbf{Aspecto} & \textbf{Microcontrolador} & \textbf{SoC con Linux} \\
\midrule
Ejemplo            & STM32, ESP32, ATmega       & RPi 5, Jetson \\
Precio             & 5--20 USD                   & 50--500 USD \\
Consumo            & 0.1--1 W                    & 3--15 W \\
RAM                & 64--512 kB                  & 1--16 GB \\
Tipo de software   & firmware bare-metal o RTOS  & sistema operativo \\
Determinismo       & alto (controlable)          & bajo (jitter del kernel) \\
Complejidad        & baja                        & alta \\
Workforce          & ingenieria embebida         & devs Linux \\
\bottomrule
\end{tabular}
\end{center}

\textbf{Cuando usar cada uno}:

\begin{itemize}[leftmargin=2em]
  \item \textbf{Microcontrolador}: para subsistemas con tareas simples y
        deterministas. EPS, ADCS, COMMS suelen usar STM32.
  \item \textbf{SoC con Linux}: para payloads que necesitan
        procesamiento pesado (IA, vision, machine learning). Tu Pi 5 es
        un caso ejemplar.
\end{itemize}

\begin{payload}
En el INTISAT, la arquitectura tipica seria:
\begin{itemize}[leftmargin=1em]
  \item \textbf{OBC principal}: microcontrolador STM32 (deterministico).
  \item \textbf{Payload}: Raspberry Pi 5 con Linux (potencia de calculo).
  \item \textbf{ADCS, COMMS, EPS}: cada uno con su propio microcontrolador.
\end{itemize}
Los buses I2C/CAN conectan todo.
\end{payload}

\chapter{Sistemas operativos}

\section{Para microcontroladores}

\textbf{Sin OS (bare-metal)}: solo loop principal + interrupciones. Util
para subsistemas muy simples.

\textbf{FreeRTOS}: el RTOS mas popular en CubeSat. Tareas, queues,
semaphores, timers. Open source.

\textbf{Zephyr}: alternativa moderna a FreeRTOS, soportado por Linux
Foundation. Tiene drivers para muchos chips.

\textbf{RTEMS}: usado en mision aeroespacial real (Mars rovers,
satellites NASA). Mas complejo de aprender.

\section{Para SoCs}

\textbf{Linux}: distribuciones tipicas en CubeSat:
\begin{itemize}[leftmargin=2em]
  \item \textbf{Raspberry Pi OS} (Bookworm): el que usamos.
  \item \textbf{Yocto / OpenEmbedded}: para builds custom super
        optimizadas.
  \item \textbf{Buildroot}: similar a Yocto pero mas simple.
  \item \textbf{Ubuntu Server}: si necesitas paquetes mas modernos.
\end{itemize}

\chapter{Frameworks de flight software}

\section{NASA cFS (Core Flight System)}

El estandar de NASA. Arquitectura publish-subscribe con apps
modulares. Usa publish/subscribe + tablas de configuracion + FDIR.

\textbf{Pros}: probado en misiones reales, ecosistema maduro.

\textbf{Contras}: complejo, curva de aprendizaje empinada.

\textbf{URL}: \url{https://github.com/nasa/cFS}

\section{NASA F' (FPrime)}

Framework moderno usado en MarCO e Ingenuity. Mas accesible que cFS.
XML define los componentes, codigo auto-generado.

\textbf{URL}: \url{https://github.com/nasa/fprime}

\section{KubOS}

Framework dedicado a CubeSat. Descontinuado pero docs todavia utiles.

\section{Custom (lo que estamos haciendo en INTISAT)}

Python + systemd corriendo en Linux. Mas simple pero menos formal.

\begin{cuidado}
Para mision \emph{academica}: custom esta bien. Para mision
\emph{comercial o critica}: usar framework formal (F' o cFS) para
trazabilidad y certificacion.
\end{cuidado}

\chapter{Tareas y concurrencia}

\section{Tareas en RTOS}

Una \emph{tarea} es un hilo de ejecucion. Cada una tiene:
\begin{itemize}[leftmargin=2em]
  \item \textbf{Prioridad}: las altas se ejecutan primero.
  \item \textbf{Stack}: memoria propia.
  \item \textbf{Estado}: running, ready, blocked, suspended.
\end{itemize}

Las tareas tipicas en un CubeSat:

\begin{center}
\begin{tabular}{lll}
\toprule
\textbf{Tarea} & \textbf{Prioridad} & \textbf{Periodo} \\
\midrule
Watchdog              & alta      & 1 s \\
Telemetria HK         & alta      & 1 s \\
Comms RX               & alta      & evento \\
Comms TX               & media     & evento \\
ADCS control           & alta      & 100 ms \\
Payload                & baja      & event-driven \\
\bottomrule
\end{tabular}
\end{center}

\section{Comunicacion inter-tarea}

\textbf{Queues}: cola FIFO de mensajes. Una tarea publica, otra
consume. Es el patron mas usado.

\textbf{Semaphores}: sincronizacion entre tareas (binarios o contadores).

\textbf{Mutexes}: proteccion de recursos compartidos.

\textbf{Notifications}: senales directas tarea-tarea (mas rapido que
queue).

\chapter{Telemetria y comandos}

\section{CCSDS Space Packet}

El estandar de empaquetado mas usado en aeroespacial:

\begin{center}
\begin{tabular}{ll}
\toprule
\textbf{Campo} & \textbf{Tamaño} \\
\midrule
Header primario             & 6 bytes \\
Header secundario           & variable \\
Datos                       & variable \\
Trailer (CRC)               & 0-4 bytes \\
\bottomrule
\end{tabular}
\end{center}

\section{Beacon}

Mensaje periodico que el satelite emite cada N segundos. Contiene
estado basico (voltaje, temperatura, modo). Es el "latido" del satelite.

\section{Housekeeping (HK)}

Empaquetado de telemetria mas detallada que el beacon. Tipicamente:
\begin{itemize}[leftmargin=2em]
  \item Voltajes de cada rail.
  \item Corrientes consumidas.
  \item Temperaturas de varios sensores.
  \item Estado de cada subsistema.
  \item Contadores de eventos.
  \item Codigos de error activos.
\end{itemize}

Periodo tipico: 1 cada 10 segundos durante operacion nominal.

\section{Comandos uplink}

Estructura tipica:
\begin{itemize}[leftmargin=2em]
  \item Opcode (que hacer).
  \item Parametros.
  \item Checksum.
  \item Authentication (cuando aplica).
\end{itemize}

\begin{payload}
Tu I2C slave del INTISAT usa una variante simplificada de CCSDS:
\texttt{[CMD\_BYTE] [LEN] [DATA...]}. Esta bien para mision academica.
Para una version mas profesional, podrias migrar a CCSDS completo.
\end{payload}

\chapter{FDIR --- Fault Detection, Isolation, Recovery}

\section{Concepto}

FDIR es el conjunto de mecanismos para detectar fallas, aislarlas, y
recuperarse autonomamente. Sin FDIR un satelite \emph{se muere} al primer
glitch.

\section{Las tres etapas}

\textbf{Deteccion}: identificar que hay un problema. Se hace con:
\begin{itemize}[leftmargin=2em]
  \item Limit checks (telemetria fuera de rango).
  \item Heartbeat checks (un subsistema no responde).
  \item Self-tests periodicos.
  \item Validacion de checksums.
\end{itemize}

\textbf{Aislamiento}: identificar QUE subsistema fallo (no propagar el
error).

\textbf{Recovery}: aplicar accion correctiva:
\begin{itemize}[leftmargin=2em]
  \item Re-intentar la operacion.
  \item Reiniciar el subsistema.
  \item Cambiar a modo de respaldo.
  \item Pasar a SAFE\_MODE.
  \item Esperar comando de tierra.
\end{itemize}

\section{Niveles de severidad}

\begin{itemize}[leftmargin=2em]
  \item \textbf{INFO}: informativo, no requiere accion.
  \item \textbf{WARNING}: prestar atencion pero no critico.
  \item \textbf{ERROR}: requiere accion.
  \item \textbf{CRITICAL}: forzar SAFE\_MODE.
\end{itemize}

\chapter{Watchdog}

\section{Watchdog hardware}

Es un timer externo al CPU. Si el CPU no le da "kick" cada N segundos,
el watchdog resetea el sistema completo.

\textbf{Por que importa}: si el CPU se cuelga (loop infinito,
deadlock, error de hardware), el watchdog es la unica forma de
recuperarse sin intervencion.

\textbf{En CubeSat}: imprescindible. Sin watchdog hardware, el satelite
se muere ante el primer cuelgue.

\section{Watchdog software (systemd)}

En Linux, systemd ofrece watchdog a nivel servicio:

\begin{lstlisting}
[Service]
Type=notify
WatchdogSec=600
ExecStart=...
\end{lstlisting}

El servicio debe llamar \texttt{sd\_notify("WATCHDOG=1")} regularmente.
Si no lo hace en 600 s, systemd lo reinicia.

\begin{payload}
Tu \texttt{cubesat-pipeline.service} ya usa WatchdogSec=600 (ver
\texttt{cubesat/systemd/*.service}). Eso protege contra cuelgues del
pipeline IA.
\end{payload}

\chapter{OTA --- Over The Air Updates}

\section{Por que se necesita}

Despues del lanzamiento no se puede mas modificar el hardware. El unico
modo de mejorar el satelite es \textbf{actualizar el software}.

\section{Requisitos del OTA}

\begin{itemize}[leftmargin=2em]
  \item \textbf{Atomico}: una actualizacion exitosa, o sino el sistema
        queda en estado previo.
  \item \textbf{Verificable}: validar el binario antes de aplicarlo.
  \item \textbf{Recuperable}: si falla, rollback automatico al estado
        previo.
  \item \textbf{Resistente a corte}: si se interrumpe (perdida de pase,
        muerte de bateria), se reintenta.
\end{itemize}

\section{Estrategias de OTA}

\textbf{A/B partitioning}: dos particiones identicas. Al actualizar,
se escribe en la inactiva. Si todo OK, se cambia el bootloader para usar
la nueva. Si falla, vuelve a la vieja.

\textbf{Git fetch + checkout}: lo que usas en el INTISAT. Simple pero
requiere disco con historico.

\textbf{Delta updates}: solo transmitir las diferencias. Ahorra bandwidth.

\begin{payload}
Tu \texttt{cubesat/ota.py} implementa git fetch + checkout con health
check + rollback. Esta validado.
\end{payload}

\chapter{Telemetria y comandos en INTISAT}

\section{Protocolo I2C}

El I2C slave en 0x42 responde a 14 comandos:

\begin{itemize}[leftmargin=2em]
  \item \texttt{CMD\_GET\_STATUS} (0x01): estado actual.
  \item \texttt{CMD\_START\_CAPTURE} (0x02): dispara scan.
  \item \texttt{CMD\_GET\_THUMBNAIL} (0x06): bajada chunked de imagen.
  \item \texttt{CMD\_GET\_TELEMETRY} (0x05): bajada de telemetria.
  \item \texttt{CMD\_SAFE\_MODE} (0x10): pasar a safe.
  \item \texttt{CMD\_RESUME} (0x11): salir de safe.
  \item \texttt{CMD\_OTA\_PREPARE} (0x20): preparar update.
  \item \texttt{CMD\_OTA\_COMMIT} (0x21): aplicar update.
  \item \texttt{CMD\_REBOOT} (0x30): reboot suave.
\end{itemize}

\section{Chunking}

El I2C limita transferencias a 32 bytes. Para data mayor (thumbnail
100 KB), se trocea en frames de 30 bytes con flag \texttt{more=1/0}.

\chapter{Configuration management}

\section{Git en orbita}

Tu daemon usa git para OTA. Eso te da:
\begin{itemize}[leftmargin=2em]
  \item Historial completo de cambios.
  \item Capacidad de rollback a cualquier commit.
  \item Branches para testing.
  \item Tags para versiones de vuelo.
\end{itemize}

\section{Tablas de parametros}

Como en cFS, se pueden cargar tablas de configuracion en runtime:

\begin{itemize}[leftmargin=2em]
  \item Limites para FDIR.
  \item Coeficientes de calibracion.
  \item Schedule de tareas.
  \item Tunable parameters del pipeline.
\end{itemize}

Esto permite ajustar comportamiento sin actualizar codigo.

\chapter{Testing del flight software}

\section{Unit tests}

Probar cada modulo aislado. Pytest, GoogleTest, etc.

\section{Integration tests}

Probar la interaccion entre modulos.

\section{HIL --- Hardware In the Loop}

Conectar el flight software a una simulacion del hardware (sensores
y actuadores virtuales). Permite probar logica completa sin satellite
real.

\section{Long duration test}

Burn-in: correr el FSW durante 168 horas (1 semana) sin reboot.
Verifica que no hay memory leaks, deadlocks acumulativos, etc.

\backmatter

\chapter*{Bibliografia}
\addcontentsline{toc}{chapter}{Bibliografia}

\begin{enumerate}[leftmargin=2em]
\item NASA. \emph{cFS documentation}. \url{https://github.com/nasa/cFS}.
\item NASA. \emph{F' Tutorials}. \url{https://nasa.github.io/fprime/}.
\item CCSDS. \emph{Space Packet Protocol Recommendation}. \url{https://public.ccsds.org}.
\item FreeRTOS. \emph{Official documentation}. \url{https://www.freertos.org}.
\item NASA Software Engineering Handbook. \url{https://swehb.nasa.gov}.
\end{enumerate}

\end{document}
